\chapter{Lesson 2: Hiragana}

\section{What hiragana is for}

Hiragana (\ja{ひらがな}) is the basic Japanese syllabary. Every Japanese word
\emph{can} be written in hiragana, and children learn it first. In ordinary adult
writing it is used for:
\begin{itemize}
  \item particles: \ja{は}, \ja{を}, \ja{に}, \ja{の}, \ldots;
  \item inflectional endings (\emph{okurigana}): the \ja{べます} of \ja{食べます};
  \item native words that have no kanji or whose kanji is rare: \ja{これ}, \ja{とても};
  \item \emph{furigana}: small reading aids printed above kanji in textbooks and
        children's books.
\end{itemize}

The shapes developed from cursive (grass-script) forms of whole kanji around the 9th
century. Hiragana was long associated with women's writing; the \emph{Tale of Genji} was
written largely in it.

\section{The basic chart (\ja{ごじゅうおん}, gojūon)}

The 46 basic hiragana are arranged in a grid called the \emph{gojūon} (``fifty sounds'').
Each row shares a consonant; each column shares a vowel. Grey cells are gaps where the
syllable either does not exist in modern Japanese or is written with another symbol.

\begin{center}
\begin{tikzpicture}[x=1.25cm,y=-1.05cm, cell/.style={draw, minimum width=1.25cm, minimum height=1.05cm, inner sep=0pt}]
  \node[font=\small\bfseries] at (1,-0.15) {a};
  \node[font=\small\bfseries] at (2,-0.15) {i};
  \node[font=\small\bfseries] at (3,-0.15) {u};
  \node[font=\small\bfseries] at (4,-0.15) {e};
  \node[font=\small\bfseries] at (5,-0.15) {o};
  \node[font=\small\bfseries, anchor=east] at (0.35,0.6) {};
  \node[cell] at (1,0.6) {};
  \node[font=\Large] at (1,0.5) {\ja{あ}};
  \node[font=\scriptsize, text=gray!50!black] at (1,0.92) {a};
  \node[cell] at (2,0.6) {};
  \node[font=\Large] at (2,0.5) {\ja{い}};
  \node[font=\scriptsize, text=gray!50!black] at (2,0.92) {i};
  \node[cell] at (3,0.6) {};
  \node[font=\Large] at (3,0.5) {\ja{う}};
  \node[font=\scriptsize, text=gray!50!black] at (3,0.92) {u};
  \node[cell] at (4,0.6) {};
  \node[font=\Large] at (4,0.5) {\ja{え}};
  \node[font=\scriptsize, text=gray!50!black] at (4,0.92) {e};
  \node[cell] at (5,0.6) {};
  \node[font=\Large] at (5,0.5) {\ja{お}};
  \node[font=\scriptsize, text=gray!50!black] at (5,0.92) {o};
  \node[font=\small\bfseries, anchor=east] at (0.35,1.6) {k};
  \node[cell] at (1,1.6) {};
  \node[font=\Large] at (1,1.5) {\ja{か}};
  \node[font=\scriptsize, text=gray!50!black] at (1,1.92) {ka};
  \node[cell] at (2,1.6) {};
  \node[font=\Large] at (2,1.5) {\ja{き}};
  \node[font=\scriptsize, text=gray!50!black] at (2,1.92) {ki};
  \node[cell] at (3,1.6) {};
  \node[font=\Large] at (3,1.5) {\ja{く}};
  \node[font=\scriptsize, text=gray!50!black] at (3,1.92) {ku};
  \node[cell] at (4,1.6) {};
  \node[font=\Large] at (4,1.5) {\ja{け}};
  \node[font=\scriptsize, text=gray!50!black] at (4,1.92) {ke};
  \node[cell] at (5,1.6) {};
  \node[font=\Large] at (5,1.5) {\ja{こ}};
  \node[font=\scriptsize, text=gray!50!black] at (5,1.92) {ko};
  \node[font=\small\bfseries, anchor=east] at (0.35,2.6) {s};
  \node[cell] at (1,2.6) {};
  \node[font=\Large] at (1,2.5) {\ja{さ}};
  \node[font=\scriptsize, text=gray!50!black] at (1,2.92) {sa};
  \node[cell] at (2,2.6) {};
  \node[font=\Large] at (2,2.5) {\ja{し}};
  \node[font=\scriptsize, text=gray!50!black] at (2,2.92) {shi};
  \node[cell] at (3,2.6) {};
  \node[font=\Large] at (3,2.5) {\ja{す}};
  \node[font=\scriptsize, text=gray!50!black] at (3,2.92) {su};
  \node[cell] at (4,2.6) {};
  \node[font=\Large] at (4,2.5) {\ja{せ}};
  \node[font=\scriptsize, text=gray!50!black] at (4,2.92) {se};
  \node[cell] at (5,2.6) {};
  \node[font=\Large] at (5,2.5) {\ja{そ}};
  \node[font=\scriptsize, text=gray!50!black] at (5,2.92) {so};
  \node[font=\small\bfseries, anchor=east] at (0.35,3.6) {t};
  \node[cell] at (1,3.6) {};
  \node[font=\Large] at (1,3.5) {\ja{た}};
  \node[font=\scriptsize, text=gray!50!black] at (1,3.92) {ta};
  \node[cell] at (2,3.6) {};
  \node[font=\Large] at (2,3.5) {\ja{ち}};
  \node[font=\scriptsize, text=gray!50!black] at (2,3.92) {chi};
  \node[cell] at (3,3.6) {};
  \node[font=\Large] at (3,3.5) {\ja{つ}};
  \node[font=\scriptsize, text=gray!50!black] at (3,3.92) {tsu};
  \node[cell] at (4,3.6) {};
  \node[font=\Large] at (4,3.5) {\ja{て}};
  \node[font=\scriptsize, text=gray!50!black] at (4,3.92) {te};
  \node[cell] at (5,3.6) {};
  \node[font=\Large] at (5,3.5) {\ja{と}};
  \node[font=\scriptsize, text=gray!50!black] at (5,3.92) {to};
  \node[font=\small\bfseries, anchor=east] at (0.35,4.6) {n};
  \node[cell] at (1,4.6) {};
  \node[font=\Large] at (1,4.5) {\ja{な}};
  \node[font=\scriptsize, text=gray!50!black] at (1,4.92) {na};
  \node[cell] at (2,4.6) {};
  \node[font=\Large] at (2,4.5) {\ja{に}};
  \node[font=\scriptsize, text=gray!50!black] at (2,4.92) {ni};
  \node[cell] at (3,4.6) {};
  \node[font=\Large] at (3,4.5) {\ja{ぬ}};
  \node[font=\scriptsize, text=gray!50!black] at (3,4.92) {nu};
  \node[cell] at (4,4.6) {};
  \node[font=\Large] at (4,4.5) {\ja{ね}};
  \node[font=\scriptsize, text=gray!50!black] at (4,4.92) {ne};
  \node[cell] at (5,4.6) {};
  \node[font=\Large] at (5,4.5) {\ja{の}};
  \node[font=\scriptsize, text=gray!50!black] at (5,4.92) {no};
  \node[font=\small\bfseries, anchor=east] at (0.35,5.6) {h};
  \node[cell] at (1,5.6) {};
  \node[font=\Large] at (1,5.5) {\ja{は}};
  \node[font=\scriptsize, text=gray!50!black] at (1,5.92) {ha};
  \node[cell] at (2,5.6) {};
  \node[font=\Large] at (2,5.5) {\ja{ひ}};
  \node[font=\scriptsize, text=gray!50!black] at (2,5.92) {hi};
  \node[cell] at (3,5.6) {};
  \node[font=\Large] at (3,5.5) {\ja{ふ}};
  \node[font=\scriptsize, text=gray!50!black] at (3,5.92) {fu};
  \node[cell] at (4,5.6) {};
  \node[font=\Large] at (4,5.5) {\ja{へ}};
  \node[font=\scriptsize, text=gray!50!black] at (4,5.92) {he};
  \node[cell] at (5,5.6) {};
  \node[font=\Large] at (5,5.5) {\ja{ほ}};
  \node[font=\scriptsize, text=gray!50!black] at (5,5.92) {ho};
  \node[font=\small\bfseries, anchor=east] at (0.35,6.6) {m};
  \node[cell] at (1,6.6) {};
  \node[font=\Large] at (1,6.5) {\ja{ま}};
  \node[font=\scriptsize, text=gray!50!black] at (1,6.92) {ma};
  \node[cell] at (2,6.6) {};
  \node[font=\Large] at (2,6.5) {\ja{み}};
  \node[font=\scriptsize, text=gray!50!black] at (2,6.92) {mi};
  \node[cell] at (3,6.6) {};
  \node[font=\Large] at (3,6.5) {\ja{む}};
  \node[font=\scriptsize, text=gray!50!black] at (3,6.92) {mu};
  \node[cell] at (4,6.6) {};
  \node[font=\Large] at (4,6.5) {\ja{め}};
  \node[font=\scriptsize, text=gray!50!black] at (4,6.92) {me};
  \node[cell] at (5,6.6) {};
  \node[font=\Large] at (5,6.5) {\ja{も}};
  \node[font=\scriptsize, text=gray!50!black] at (5,6.92) {mo};
  \node[font=\small\bfseries, anchor=east] at (0.35,7.6) {y};
  \node[cell] at (1,7.6) {};
  \node[font=\Large] at (1,7.5) {\ja{や}};
  \node[font=\scriptsize, text=gray!50!black] at (1,7.92) {ya};
  \node[cell, fill=gray!15] at (2,7.6) {};
  \node[cell] at (3,7.6) {};
  \node[font=\Large] at (3,7.5) {\ja{ゆ}};
  \node[font=\scriptsize, text=gray!50!black] at (3,7.92) {yu};
  \node[cell, fill=gray!15] at (4,7.6) {};
  \node[cell] at (5,7.6) {};
  \node[font=\Large] at (5,7.5) {\ja{よ}};
  \node[font=\scriptsize, text=gray!50!black] at (5,7.92) {yo};
  \node[font=\small\bfseries, anchor=east] at (0.35,8.6) {r};
  \node[cell] at (1,8.6) {};
  \node[font=\Large] at (1,8.5) {\ja{ら}};
  \node[font=\scriptsize, text=gray!50!black] at (1,8.92) {ra};
  \node[cell] at (2,8.6) {};
  \node[font=\Large] at (2,8.5) {\ja{り}};
  \node[font=\scriptsize, text=gray!50!black] at (2,8.92) {ri};
  \node[cell] at (3,8.6) {};
  \node[font=\Large] at (3,8.5) {\ja{る}};
  \node[font=\scriptsize, text=gray!50!black] at (3,8.92) {ru};
  \node[cell] at (4,8.6) {};
  \node[font=\Large] at (4,8.5) {\ja{れ}};
  \node[font=\scriptsize, text=gray!50!black] at (4,8.92) {re};
  \node[cell] at (5,8.6) {};
  \node[font=\Large] at (5,8.5) {\ja{ろ}};
  \node[font=\scriptsize, text=gray!50!black] at (5,8.92) {ro};
  \node[font=\small\bfseries, anchor=east] at (0.35,9.6) {w};
  \node[cell] at (1,9.6) {};
  \node[font=\Large] at (1,9.5) {\ja{わ}};
  \node[font=\scriptsize, text=gray!50!black] at (1,9.92) {wa};
  \node[cell, fill=gray!15] at (2,9.6) {};
  \node[cell, fill=gray!15] at (3,9.6) {};
  \node[cell, fill=gray!15] at (4,9.6) {};
  \node[cell] at (5,9.6) {};
  \node[font=\Large] at (5,9.5) {\ja{を}};
  \node[font=\scriptsize, text=gray!50!black] at (5,9.92) {o};
  \node[font=\small\bfseries, anchor=east] at (0.35,10.6) {};
  \node[cell] at (1,10.6) {};
  \node[font=\Large] at (1,10.5) {\ja{ん}};
  \node[font=\scriptsize, text=gray!50!black] at (1,10.92) {n};
  \node[cell, fill=gray!15] at (2,10.6) {};
  \node[cell, fill=gray!15] at (3,10.6) {};
  \node[cell, fill=gray!15] at (4,10.6) {};
  \node[cell, fill=gray!15] at (5,10.6) {};
\end{tikzpicture}
\end{center}

\begin{important}[Irregular romanizations in the chart]
Most cells are just consonant + vowel, but note:
\ja{し} shi (not si), \ja{ち} chi (not ti), \ja{つ} tsu (not tu), \ja{ふ} fu (not hu).
\ja{を} is pronounced ``o'' and is used \emph{only} as the object particle.
\ja{ん} is the only kana that is a consonant alone.
\end{important}

\subsection{Stroke order and look-alikes}

General stroke-order rules (shared with kanji): top to bottom, left to right, horizontal
before a crossing vertical. Hiragana are written with curved strokes; the end of each
stroke is a ``stop'' (\emph{tome}), ``flick'' (\emph{hane}) or ``sweep'' (\emph{harai}),
and the quiz markers do check these.

Pairs and groups that beginners confuse:
\begin{center}
\begin{tabular}{ll}
\toprule
Group & How to tell them apart \\
\midrule
\ja{あ} a / \ja{お} o & \ja{お} has a separate dot at upper right \\
\ja{い} i / \ja{り} ri & \ja{り} is taller; left stroke hooks up into the right \\
\ja{さ} sa / \ja{ち} chi & mirror images: \ja{さ} opens left at the bottom, \ja{ち} opens right \\
\ja{ぬ} nu / \ja{め} me & \ja{ぬ} has a small loop at the end \\
\ja{ね} ne / \ja{れ} re / \ja{わ} wa & loop (ne), outward flick (re), inward curve (wa) \\
\ja{は} ha / \ja{ほ} ho & \ja{ほ} has an extra horizontal stroke on top \\
\ja{る} ru / \ja{ろ} ro & \ja{る} ends in a loop \\
\ja{く} ku / \ja{へ} he & \ja{く} opens right (``<''), \ja{へ} is a low hill \\
\bottomrule
\end{tabular}
\end{center}

\begin{aside}
Mnemonics that worked for me: \ja{く} is a beak going ``\emph{cuckoo}''; \ja{ぬ} is a
bowl of \emph{noodles} with chopsticks; \ja{ね} is a \emph{cat} (\ja{ねこ}) curled up
with its tail; \ja{つ} is a \emph{tsunami} wave; \ja{し} is a \emph{she}ep's crook. Make
up your own: the silly ones stick best.
\end{aside}

\section{Dakuten and handakuten (voiced sounds)}

Two small marks change a kana's consonant:
\begin{itemize}
  \item \textbf{dakuten} \ja{゛} (two short strokes, nicknamed \emph{ten-ten}) voices the
        consonant: k$\to$g, s$\to$z, t$\to$d, h$\to$b.
  \item \textbf{handakuten} \ja{゜} (a small circle, \emph{maru}) turns h into p.
\end{itemize}

\begin{center}
\begin{tikzpicture}[box/.style={draw, rounded corners, minimum width=1.3cm, minimum height=1cm, font=\Large}]
  \node[box] (k) at (0,0) {\ja{か}};
  \node[box, fill=blue!8] (g) at (3,0) {\ja{が}};
  \draw[-Stealth, thick] (k) -- node[above, font=\small] {+ \ja{゛}} (g);
  \node[font=\small] at (0,-0.8) {ka}; \node[font=\small] at (3,-0.8) {ga};
  \node[box] (h) at (7,0) {\ja{は}};
  \node[box, fill=blue!8] (b) at (10,0.9) {\ja{ば}};
  \node[box, fill=red!8] (p) at (10,-0.9) {\ja{ぱ}};
  \draw[-Stealth, thick] (h) -- node[above left, font=\small] {+ \ja{゛}} (b);
  \draw[-Stealth, thick] (h) -- node[below left, font=\small] {+ \ja{゜}} (p);
  \node[font=\small] at (7,-0.8) {ha};
  \node[font=\small, anchor=west] at (10.8,0.9) {ba};
  \node[font=\small, anchor=west] at (10.8,-0.9) {pa};
\end{tikzpicture}
\end{center}

\begin{center}
\begin{tabular}{lccccc}
\toprule
 & a & i & u & e & o \\
\midrule
g & \ja{が} ga & \ja{ぎ} gi & \ja{ぐ} gu & \ja{げ} ge & \ja{ご} go \\
z & \ja{ざ} za & \ja{じ} ji & \ja{ず} zu & \ja{ぜ} ze & \ja{ぞ} zo \\
d & \ja{だ} da & \ja{ぢ} ji & \ja{づ} zu & \ja{で} de & \ja{ど} do \\
b & \ja{ば} ba & \ja{び} bi & \ja{ぶ} bu & \ja{べ} be & \ja{ぼ} bo \\
p & \ja{ぱ} pa & \ja{ぴ} pi & \ja{ぷ} pu & \ja{ぺ} pe & \ja{ぽ} po \\
\bottomrule
\end{tabular}
\end{center}

\begin{remark}
\ja{ぢ} and \ja{づ} sound the same as \ja{じ} and \ja{ず}. They are rare and appear mainly
when a word's \ja{ち}/\ja{つ} gets voiced in a compound or repetition:
\ja{はな} + \ja{ち} $\to$ \ja{はなぢ} \emph{hanaji} ``nosebleed'',
\ja{つづく} \emph{tsuzuku} ``to continue''. Otherwise always use \ja{じ} and \ja{ず}.
\end{remark}

Vocabulary to practise reading:
\begin{center}
\begin{tabular}{lll}
\toprule
Hiragana & Romaji & Meaning \\
\midrule
\ja{かぎ} & kagi & key \\
\ja{ざっし} & zasshi & magazine \\
\ja{でんわ} & denwa & telephone \\
\ja{ぶどう} & budō & grapes \\
\ja{えんぴつ} & enpitsu & pencil \\
\ja{たんぽぽ} & tanpopo & dandelion \\
\bottomrule
\end{tabular}
\end{center}

\section{Contracted sounds: small \ja{ゃ}, \ja{ゅ}, \ja{ょ} (yōon)}

A kana from the \emph{i}-column followed by a \emph{small} \ja{ゃ}, \ja{ゅ} or \ja{ょ}
makes one combined mora: \ja{き} + \ja{ゃ} = \ja{きゃ} kya.
\begin{center}
\begin{tabular}{lccc}
\toprule
 & ya & yu & yo \\
\midrule
k & \ja{きゃ} kya & \ja{きゅ} kyu & \ja{きょ} kyo \\
s & \ja{しゃ} sha & \ja{しゅ} shu & \ja{しょ} sho \\
t & \ja{ちゃ} cha & \ja{ちゅ} chu & \ja{ちょ} cho \\
n & \ja{にゃ} nya & \ja{にゅ} nyu & \ja{にょ} nyo \\
h & \ja{ひゃ} hya & \ja{ひゅ} hyu & \ja{ひょ} hyo \\
m & \ja{みゃ} mya & \ja{みゅ} myu & \ja{みょ} myo \\
r & \ja{りゃ} rya & \ja{りゅ} ryu & \ja{りょ} ryo \\
g & \ja{ぎゃ} gya & \ja{ぎゅ} gyu & \ja{ぎょ} gyo \\
j & \ja{じゃ} ja & \ja{じゅ} ju & \ja{じょ} jo \\
b & \ja{びゃ} bya & \ja{びゅ} byu & \ja{びょ} byo \\
p & \ja{ぴゃ} pya & \ja{ぴゅ} pyu & \ja{ぴょ} pyo \\
\bottomrule
\end{tabular}
\end{center}

\begin{important}[Size matters]
The small kana is written about half size, in the lower-left of its square (in
horizontal writing). Compare:
\ja{びょういん} \emph{byōin} ``hospital'' (4 morae) vs.\
\ja{びよういん} \emph{biyōin} ``beauty salon'' (5 morae).
\end{important}

\section{The small \ja{っ}: doubled consonants (sokuon)}

A small \ja{っ} before a consonant doubles that consonant. It is pronounced as a
one-mora pause (silence with the mouth ready for the next consonant). In romaji, double
the next consonant letter.
\begin{itemize}
  \item \ja{きって} \emph{kitte} ``postage stamp'' vs.\ \ja{きて} \emph{kite} ``come (te-form)''
  \item \ja{がっこう} \emph{gakkō} ``school''
  \item \ja{ざっし} \emph{zasshi} ``magazine''
  \item \ja{いっぱい} \emph{ippai} ``full''
  \item \ja{まっちゃ} \emph{matcha} ``powdered green tea'' (before \emph{ch} write \emph{tch})
\end{itemize}
\ja{っ} never begins a word and does not normally appear before vowels, n, m, r, w
or y in standard spelling (in exclamations like \ja{あっ！} ``Ah!'' it marks a cut-off).

\section{Long vowels in hiragana}

\begin{center}
\begin{tabular}{llll}
\toprule
Long vowel & Written as & Example & Meaning \\
\midrule
ā & + \ja{あ} & \ja{おかあさん} okāsan & mother \\
ī & + \ja{い} & \ja{おにいさん} onīsan & older brother \\
ū & + \ja{う} & \ja{すうがく} sūgaku & mathematics \\
ē & + \ja{い} (usually) & \ja{せんせい} sensei & teacher \\
ē & + \ja{え} (a few words) & \ja{おねえさん} onēsan & older sister \\
ō & + \ja{う} (usually) & \ja{ひこうき} hikōki & airplane \\
ō & + \ja{お} (a few words) & \ja{おおきい} ōkii & big \\
\bottomrule
\end{tabular}
\end{center}

\begin{remark}
\ja{えい} and \ja{おう} are written that way for historical reasons but in normal speech
are just long \emph{ē} and \emph{ō}. The \ja{おお} spelling words are few enough to
memorise: \ja{おおきい} big, \ja{とお} ten (things), \ja{とおい} far, \ja{おおい} many.
\end{remark}

\section{Particles that break the spelling rules}
Three particles keep old spellings:
\begin{center}
\begin{tabular}{llll}
\toprule
Particle & Written & Pronounced & Function \\
\midrule
topic & \ja{は} & wa & ``as for \ldots'' \\
direction & \ja{へ} & e & ``toward \ldots'' \\
object & \ja{を} & o & marks the direct object \\
\bottomrule
\end{tabular}
\end{center}
Example: \ja{わたしはがっこうへいきます。} \emph{Watashi wa gakkō e ikimasu.}
``I go to school.''

\section{Punctuation and writing direction}
\begin{itemize}
  \item Full stop \ja{。} (\emph{maru}), comma \ja{、} (\emph{ten}), quotation brackets
        \ja{「」}.
  \item No spaces between words in normal writing (beginner textbooks add them).
  \item Text can be written horizontally (left to right) or vertically (top to bottom,
        columns right to left). Novels and newspapers are mostly vertical.
  \item Each character fills an imaginary square of the same size; this is why grid
        paper (\ja{げんこうようし}, \emph{genkō yōshi}) is used for compositions.
\end{itemize}

\section*{Lesson 2 review}
\begin{itemize}
  \item 46 basic hiragana in the gojūon grid (rows k, s, t, n, h, m, y, r, w + \ja{ん}).
  \item \ja{゛} voices (g, z, d, b); \ja{゜} gives p.
  \item Small \ja{ゃゅょ} after an \emph{i}-column kana: one mora (\ja{きょ} kyo).
  \item Small \ja{っ}: one-mora pause, doubles the next consonant (\ja{きって} kitte).
  \item Long \emph{ē} and \emph{ō} are usually spelled \ja{えい}, \ja{おう}.
  \item Particles \ja{は}, \ja{へ}, \ja{を} are read wa, e, o.
\end{itemize}
